\section{Interpretation and limits}
\label{sec:discussion}
The measured gain has several sources. Exact identities remove an elimination problem when their assumptions hold. Finite-field reconstruction changes how Kira carries out the remaining coefficient algebra. Shared calls reduce repeated initialization, and learned choices can reduce the equation system in domains where a matched test supports them. The separate controls in Section~\ref{sec:results} show why the combined timing improvement cannot all be attributed to machine learning.

The physically relevant distinction is the structure of the propagators. An independent thirteen-loop product can be reduced by multiplying known component results, and a connected bubble tree can be integrated from its leaves. A graph with additional momentum couplings need not permit either step. Table~\ref{tab:exact-domains} and the appendix derivations state the domains of the formulas, while Table~\ref{tab:coverage} states the smaller set that received the final complete-time test.

The native benchmarks keep $d$ symbolic and fix the masses and scattering invariants to rational values. Their outputs are exact rational functions of $d$. The analytic scalar mass formulas also admit symbolic squared masses in their stated domain. These tests do not establish complete-time performance for arbitrary symbolic multiscale scattering, all numerator powers, or every connected graph through thirteen loops. A finite comparison also cannot prove that the chosen settings minimize time globally.

Recovery deserves particular attention because the cost of a failed restricted system can exceed the work saved by removing equations. The latest equation-filter trial exhibits this effect directly. Its final outputs agree with their references after recovery, but its time is 71.1\% greater. Predicting failure from training data can guide a choice, while completion still requires algebraic checks. Exact formulas have a different admission test: their propagator and kinematic conditions must be verified before substitution.

The final application thus supports a concrete use case: requests drawn from the stated analytic domains and native families, with the tested basis and arithmetic settings. Its 8.8\% aggregate gain is accompanied by slower subgroups, and the latest model-only evidence remains inconclusive or negative. Wider performance claims require new independent graph structures and complete symbolic-kinematic comparisons with equally tuned controls.
